% Real matched case: 09709b691b1e860345ead2f5, seed17.
% Source: qualitative_cases_abc.json; lines are retrieved facts, not attention traces.
\resizebox{\linewidth}{!}{%
\begin{tikzpicture}[x=1.35cm,y=1cm,font=\fontsize{9.5}{11.5}\selectfont,
 n/.style={draw=gray!65,rounded corners=1pt,align=center,minimum height=.55cm,minimum width=1.7cm,inner sep=3pt},
 edge/.style={->,>=stealth,blue!65!black,line width=.7pt},
 absent/.style={edge,red!70!black,dashed},
 lab/.style={fill=white,inner sep=1pt,font=\fontsize{9}{11}\selectfont},
 toplab/.style={lab,midway}]
\node[anchor=west,align=left,text width=18.3cm] at (0,15.3) {问题：除A-5691424之外，同城同街道的其他事故明确记录了哪些受影响部位？\\标准答案：右路肩（right shoulder）。};
\begin{scope}[yshift=13.5cm]
\node[anchor=west,font=\fontsize{10}{12}\selectfont\bfseries] at (0,.95) {(a) 去除检索中的问题更新};
\node[n] (aroot) at (1,0) {A-5691424};
\node[n] (aroad) at (3.8,0) {I-95};
\node[n] (atarget) at (6.6,0) {A-5396309};
\node[n] (atb) at (9.5,0) {目标阻断事件};
\node[n,] (aright) at (12.6,0) {右路肩};
\node[n] (asb) at (3.8,-1.35) {原事故阻断事件};
\node[n,draw=orange!85!black,line width=1pt] (aleft) at (7.2,-1.35) {左车道};
\node[n,] (astate) at (11,-1.35) {blocked（阻断）};
\draw[edge] (aroot)--node[toplab]{道路}(aroad);
\draw[edge] (atarget)--node[toplab]{道路}(aroad);
\draw[absent] (atarget)--node[toplab]{包含事件}(atb);
\draw[absent] (atb)--node[toplab]{受影响部位}(aright);
\draw[edge] (aroot)--node[toplab,sloped]{包含事件}(asb);
\draw[edge] (asb)--node[toplab]{受影响部位}(aleft);
\draw[edge] (asb.south)--++(0,-.85)--node[toplab]{阻断状态}(11,-2.475)--(astate.south);
\node[anchor=west] at (0,-3.2) {候选前3名：左车道；blocked；其他阻断事件。支持事实覆盖：2/4。};
\end{scope}
\begin{scope}[yshift=8.8cm]
\node[anchor=west,font=\fontsize{10}{12}\selectfont\bfseries] at (0,.95) {(b) 去除Mamba选择性过滤};
\node[n] (broot) at (1,0) {A-5691424};
\node[n] (broad) at (3.8,0) {I-95};
\node[n] (btarget) at (6.6,0) {A-5396309};
\node[n] (btb) at (9.5,0) {目标阻断事件};
\node[n,] (bright) at (12.6,0) {右路肩};
\node[n] (bsb) at (3.8,-1.35) {原事故阻断事件};
\node[n,] (bleft) at (7.2,-1.35) {左车道};
\node[n,draw=orange!85!black,line width=1pt] (bstate) at (11,-1.35) {blocked（阻断）};
\draw[edge] (broot)--node[toplab]{道路}(broad);
\draw[edge] (btarget)--node[toplab]{道路}(broad);
\draw[edge] (btarget)--node[toplab]{包含事件}(btb);
\draw[edge] (btb)--node[toplab]{受影响部位}(bright);
\draw[edge] (broot)--node[toplab,sloped]{包含事件}(bsb);
\draw[edge] (bsb)--node[toplab]{受影响部位}(bleft);
\draw[edge] (bsb.south)--++(0,-.85)--node[toplab]{阻断状态}(11,-2.475)--(bstate.south);
\node[anchor=west] at (0,-3.2) {候选前3名：blocked；右路肩；左车道。支持事实覆盖：4/4。};
\end{scope}
\begin{scope}[yshift=4.1cm]
\node[anchor=west,font=\fontsize{10}{12}\selectfont\bfseries] at (0,.95) {(c) 完整模型};
\node[n] (croot) at (1,0) {A-5691424};
\node[n] (croad) at (3.8,0) {I-95};
\node[n] (ctarget) at (6.6,0) {A-5396309};
\node[n] (ctb) at (9.5,0) {目标阻断事件};
\node[n,draw=orange!85!black,line width=1pt] (cright) at (12.6,0) {右路肩};
\node[n] (csb) at (3.8,-1.35) {原事故阻断事件};
\node[n,] (cleft) at (7.2,-1.35) {左车道};
\node[n,] (cstate) at (11,-1.35) {blocked（阻断）};
\draw[edge] (croot)--node[toplab]{道路}(croad);
\draw[edge] (ctarget)--node[toplab]{道路}(croad);
\draw[edge] (ctarget)--node[toplab]{包含事件}(ctb);
\draw[edge] (ctb)--node[toplab]{受影响部位}(cright);
\draw[edge] (croot)--node[toplab,sloped]{包含事件}(csb);
\draw[edge] (csb)--node[toplab]{受影响部位}(cleft);
\draw[edge] (csb.south)--++(0,-.85)--node[toplab]{阻断状态}(11,-2.475)--(cstate.south);
\node[anchor=west] at (0,-3.2) {候选前3名：右路肩；左车道；blocked。支持事实覆盖：4/4。};
\end{scope}
\node[anchor=west,font=\fontsize{9.5}{11.5}\selectfont] at (0,.15) {蓝色实线：已检索关系；红色虚线：缺失的标准支持关系；橙色边框：预测答案。};
\end{tikzpicture}
}
